\subsection*{NT-A08--A11. Prescribed ray class and Hecke angle, with uniform parameters}

Let \(K\) have degree \(n\) and absolute discriminant \(D\).
For a primitive unitary Hecke character \(\lambda\), use the standard
finite-part completion, with the complex factors split by duplication:
\[
 \Lambda_K(s,\lambda)=
 q_\lambda^{s/2}\prod_{j=1}^{n}\Gamma_{\mathbf R}(s+\kappa_{\lambda,j})
 L_K(s,\lambda),\qquad
 q_\lambda=D\,N\mathfrak f_\lambda,\qquad
 \Gamma_{\mathbf R}(s)=\pi^{-s/2}\Gamma(s/2).
                                                        \tag{A08.1}
\]
An irrelevant constant of absolute value one in the functional equation
is allowed. At a real place the parameter is
\(\epsilon+i\tau\), \(\epsilon\in\{0,1\}\); at a complex place the
two parameters are \(|k|/2+i\tau\) and \(|k|/2+i\tau+1\).
Put
\[
 {\cal C}(\lambda)=q_\lambda
                    \prod_{j=1}^{n}(3+|\kappa_{\lambda,j}|).
\]
These are the local Hecke conventions of Ramakrishnan--Valenza
Theorem~7-19 and its local-factor discussion, also checked in
MT-M28; Rajan, \S1, pp.~281--283, records the conductor and
infinity-type behavior under powers.

\paragraph{The exact classical prime-sum input.}
Iwaniec--Kowalski, \emph{Analytic Number Theory}, Theorem~5.15,
printed p.~114, equation~(5.56), states, under RH and the
Ramanujan bound for a function of degree \(d\), that
\[
 \psi_f(x)=r_fx+
 O\bigl(\sqrt x\log(2x)
             [\log{\cal C}(f)+d\log(2x)]\bigr),\qquad x\ge2,
                                                        \tag{A08.2}
\]
with an absolute constant. Here \(\psi_f\) uses the coefficients
of \(-L_f'/L_f\), \(r_f\) is the pole order at \(1\), and
\({\cal C}(f)\) is the analytic conductor at zero, defined in
\S5.1, (5.6)--(5.7), p.~95.
Its preceding proof sentence uses the truncated explicit formula
(5.53), pp.~112--113, and the local zero count (5.26).
This is a cited classical theorem; its preceding exercise and
full analytic development are not reproved here.

The book's definition in \S5.1, pp.~93--95, requires its infinity
parameters to be real or occur in conjugate pairs. An arbitrary
individual unitary Hecke function need not meet that convention.
The following complete reduction removes that restriction without
enlarging the order of the error term.

\paragraph{NT-A08: a uniform complex Hecke prime-sum bound.}
Suppose \(\lambda\) is not a pure unitary norm character and
RH holds for \(L_K(s,\lambda)\). Then, with an absolute constant,
\[
 \left|\sum_{N\mathfrak a\le x}
       \Lambda_K(\mathfrak a)\lambda(\mathfrak a)\right|
 \ll \sqrt x\log(2x)
          [\log{\cal C}(\lambda)+n\log(2x)].             \tag{A08.3}
\]
The value at an ideal not coprime to the conductor is zero.
Consequently, for any integral ideal \(\mathfrak f\) divisible by
\(\mathfrak f_\lambda\),
\[
 \sum_{\substack{N\mathfrak p\le x\\\mathfrak p\nmid\mathfrak f}}
            \lambda(\mathfrak p)
 \ll \sqrt x[\log{\cal C}(\lambda)+n\log(2x)]
                       +\log N\mathfrak f.             \tag{A08.4}
\]
For \(\lambda=1\), the same bound applies after subtracting
\(\operatorname{Li}_2(x)\) from the left side, using RH for
\(\zeta_K\). This principal case is handled separately.

\begin{proof}
Write
\[
 -\frac{L'}L(s,\lambda)=\sum_{r\ge1}b_r r^{-s}.
 \]
Grouping the ideal Euler factors by rational prime gives
\(|b_r|\le n\Lambda(r)\). In fact at \(r=p^a\) the coefficient
is a sum of terms \(f_{\mathfrak p}\log p\) times unit complex
numbers, over residue degrees \(f_{\mathfrak p}\mid a\), and
\(\sum_{\mathfrak p\mid p}f_{\mathfrak p}\le n\).
The local roots of the resulting degree-\(n\) rational-prime
Euler product have absolute value at most one.

For real \(t\) form
\[
 F_t(s)=L_K(s+it,\lambda)L_K(s-it,\overline\lambda).
 \]
Because \(\lambda\) is not pure norm, both factors are entire;
there are no shifted poles. Its degree is \(2n\), its conductor
is \(q_\lambda^2\), and its infinity parameters are
\(\kappa_{\lambda,j}+it\) and their conjugates. Its functional
equation has root number one. These statements include the
bad-prime Euler factors. The Hecke continuation, finite order
and boundary nonvanishing supply the remaining axioms of
Iwaniec--Kowalski. RH for the two translated factors follows
from RH for \(\lambda\) and complex conjugation.
For \(|t|\le2\),
\[
 \log{\cal C}(F_t)\le
           2\log{\cal C}(\lambda)+2n\log2 .
 \]
Its logarithmic derivative has coefficient
\(2\operatorname{Re}(b_r r^{-it})\).
Thus (A08.2), with pole order zero, bounds
\[
 R_t(y):=\sum_{r\le y}\operatorname{Re}(b_r r^{-it}),
 \qquad
 |R_t(y)|\le C\sqrt y\log(2y)
           [\log{\cal C}(\lambda)+n\log(2y)]              \tag{A08.5}
\]
uniformly in \(|t|\le2\).

Here are both steps recovering the imaginary part. Fix \(x\ge4\),
put \(a=\sqrt x\) and \(t=\pi/(2\log x)\), and hold this \(t\)
fixed while \(u\) ranges over \([a,x]\).
Denote the right side of (A08.5), with \(y=x\), by \(B_x\).
Write \(b_r=u_r+iv_r\). First partial summation gives
\[
 \begin{split}
 \sum_{a<r\le y}u_r\cos(t\log r)
  ={}&[R_0(y)-R_0(a)]\cos(t\log y)\\
    &+\int_a^y[R_0(u)-R_0(a)]
                       \frac{t\sin(t\log u)}u\,du .
 \end{split}
\]
The total variation of the cosine on this interval is at most
one, so this sum is \(O(B_x)\), uniformly for \(a\le y\le x\).
Since
\(\operatorname{Re}(b_r r^{-it})
   =u_r\cos(t\log r)+v_r\sin(t\log r)\),
it follows that
\[
 V(y):=\sum_{a<r\le y}v_r\sin(t\log r)=O(B_x).
 \]
Throughout \([a,x]\) the sine is at least \(1/\sqrt2\).
A second partial summation gives
\[
 \sum_{a<r\le x}v_r
   =\frac{V(x)}{\sin(t\log x)}
    +\int_a^x V(u)
                \frac{t\cos(t\log u)}
                     {u\sin^2(t\log u)}\,du
   =O(B_x),
 \]
because \(1/\sin(t\log u)\) has total variation at most
\(\sqrt2-1\). The initial segment has absolute sum at most
\(n\psi(\sqrt x)\ll n\sqrt x\), by the elementary Chebyshev
bound. It is absorbed in \(B_x\), since
\(\log{\cal C}(\lambda)\ge n\log3\).
The real part was already bounded by \(R_0(x)\).
This proves (A08.3); \(2\le x<4\) follows directly from the
coefficient bound.

Removing higher prime powers costs \(O(n\sqrt x\log(2x))\),
which is absorbed in (A08.3). Indeed, their underlying rational
primes are at most \(\sqrt x\), and the total weight above
each is at most \(n\log x\).
Abel summation removes the prime logarithm: if the weighted
error is \(O(\sqrt u\log(2u)E(u))\), with
\(E(u)=\log{\cal C}(\lambda)+n\log(2u)\), the endpoint error
is \(O(\sqrt xE(x))\), and its integral is bounded by
\[
 E(x)\int_2^x\frac{\log(2u)}
                     {\sqrt u(\log u)^2}\,du
                       \ll\sqrt x E(x).
 \]
Omitting additionally the primes dividing \(\mathfrak f\)
costs at most \(\omega(\mathfrak f)\le
\log N\mathfrak f/\log2\).
This proves (A08.4). For the principal character use the
Dedekind prime-count estimate NT-C02, with the same deletion.
In particular the principal factor is never put into \(F_t\).
\end{proof}

\paragraph{NT-A09: joint class and arc discrepancy.}
Let \(\mathfrak m=\mathfrak m_0\mathfrak m_\infty\) be a ray
modulus, \(G=\operatorname{Cl}_{\mathfrak m}(K)\) and \(h=|G|\).
Let \(\eta\) be unitary and not of the form finite order times
a unitary norm power. Put
\[
 \begin{aligned}
 \mathfrak f&=\operatorname{lcm}(\mathfrak m_0,\mathfrak f_\eta),
 & B_\eta&=4+\max_j|\kappa_{\eta,j}|,\\
 Q&=D\,N\mathfrak f\,B_\eta^n,
 & {\cal L}(x)&=\log Q+2n\log(3x).
 \end{aligned}                                                     \tag{A09.1}
\]
For \(a\in G\) and an arc \(J\subset\mathbf R/\mathbf Z\), let
\(N_{a,J}(x)\) count primes \(\mathfrak p\nmid\mathfrak f\)
of norm at most \(x\) with ray class \(a\) and
\(\eta(\mathfrak p)\in e^{2\pi iJ}\).
Assume RH for the primitive Hecke functions attached to
\[
             \chi\eta^m,\qquad
     \chi\in\widehat G,\quad m=0,1,2,\ldots .             \tag{A09.2}
\]
Conditional input \(M\) suffices. Then
\[
 N_{a,J}(x)=\frac{|J|}{h}\operatorname{Li}_2(x)
           +O\bigl(\sqrt x\log(2x){\cal L}(x)\bigr),
 \qquad x\ge2,                                         \tag{A09.3}
\]
with an absolute constant, uniformly in \(K,\eta,\mathfrak m,a,J\).
The same bound holds if only absolute-degree-one primes are counted.
In particular there is no factor \(h\) in the error.

\begin{proof}
For \(\chi\in\widehat G\), its finite conductor divides
\(\mathfrak m_0\); its infinity type has only real signs.
The finite conductor of \(\chi\eta^m\) divides \(\mathfrak f\).
At complex places its angular and imaginary parameters grow
by \(m\); at real places its parity remains in \(\{0,1\}\).
The conventions (A08.1) therefore give directly
\[
 {\cal C}(\chi\eta^m)\le Q(m+1)^n,\qquad m\ge0.           \tag{A09.4}
\]
For \(m>0\), none of these characters is pure norm. Otherwise,
if \(\chi\eta^m=|\cdot|^{it}\), the character
\(\eta|\cdot|^{-it/m}\) would have finite order: its \(m\)-th
power is the finite-order character \(\chi^{-1}\).
This contradicts the hypothesis on \(\eta\).
For \(m=0\), only \(\chi=1\) is principal.

Apply (A08.4), treating that principal character separately.
The additional \(\log N\mathfrak f\) is absorbed by
\(\sqrt x\log Q\). For each \(m\ge0\) we obtain
\[
 \sum_{\substack{N\mathfrak p\le x\\\mathfrak p\nmid\mathfrak f}}
       \chi(\mathfrak p)\eta(\mathfrak p)^m
 =\mathbf1_{\chi=1,m=0}\operatorname{Li}_2(x)
 +O\left(\sqrt x\{\log Q+n\log[(m+1)(2x)]\}\right).
                                                               \tag{A09.5}
\]
The ray-class indicator is
\[
 \mathbf1_{[\mathfrak p]=a}
       =\frac1h\sum_{\chi\in\widehat G}
                        \overline{\chi(a)}\chi(\mathfrak p).
 \]
Its \(h\) terms cancel the prefactor \(1/h\) when the errors
in (A09.5) are summed in absolute value. Thus the number
\(N_a(x)\) in class \(a\) has main term \(\operatorname{Li}_2(x)/h\),
and every nonzero angular Fourier sum in that class has
the same error bound without \(h\).

Apply the Erd\H{o}s--Tur\'an inequality S-NT-ET to the finite
list of angles in this single class, with \(H=\lceil x\rceil\).
Its zero-frequency truncation error is
\(O(N_a(x)/(H+1))=O(n)\): at most \(n\) prime ideals lie over
each rational prime of size at most \(x\).
For \(1\le m\le H\),
\(\log Q+n\log[(m+1)(2x)]\ll{\cal L}(x)\).
Summing the \(1/m\) weights gives (A09.3), first for half-open
arcs. Singleton bounds follow by decreasing arcs to an endpoint;
changing either endpoint convention changes the count by at
most the same error. This also covers arcs crossing zero.
Finally there are at most \(n\pi(\sqrt x)=O(n\sqrt x)\)
prime ideals of absolute residue degree at least two and norm
at most \(x\). Their deletion is absorbed in the error.
\end{proof}

\paragraph{NT-A10: shrinking arcs with separate class-number dependence.}
Fix \(c>0\) and \(0<\alpha<1/2\), and let
\(u_0=\max\{2,c^{1/\alpha}\}\).
For a fixed center let the arc at norm \(u>u_0\) have length
\(c u^{-\alpha}\). Let \(N_{a,c,\alpha}(x)\) count precisely
the primes \(\mathfrak p\nmid\mathfrak f\) with
\(u_0<N\mathfrak p\le x\), ray class \(a\), and angle in
that arc at \(u=N\mathfrak p\).
For \(x>\max\{4,u_0\}\), this joint count satisfies
\[
 \begin{split}
 N_{a,c,\alpha}(x)
  ={}&\frac c h\int_{u_0}^x\frac{u^{-\alpha}}{\log u}\,du\\
   &+O_{\alpha,c}\left(
       \sqrt x\log(2x){\cal L}(x)
       +h^{-1/2}x^{3/4-\alpha/2}\sqrt{{\cal L}(x)}
                  \right).                            
 \end{split} \tag{A10.1}
\]
The constant is uniform in the center and in the parameters
of (A09.1). For varying fields and ray moduli, a sufficient
uniform asymptotic regime is
\[
 h\,{\cal L}(x)(\log x)^2=o(x^{1/2-\alpha}).              \tag{A10.2}
\]
For fixed data this holds automatically. The same assertions
hold after restricting to degree-one primes.

\begin{proof}
Set \(T(x)=\log(2x){\cal L}(x)\). Partition \((u_0,x]\)
geometrically with ratio \(1+\delta\), \(0<\delta\le1/2\).
Nested centered arcs bound the moving count on each interval
between its two fixed-endpoint arc counts. Apply (A09.3) to
each difference of prime-counting functions. The geometric
sum of endpoint errors is
\[
 O\bigl(\delta^{-1}\sqrt x\,T(x)\bigr),
 \]
since \(\sum \sqrt{u_j}\ll\delta^{-1}\sqrt x\), and
\(T(u_j)\le T(x)\).
The two main terms bracket the displayed integral, and their
difference from it is
\[
 O_\alpha\left(\frac c h\delta
           \int_{u_0}^x\frac{u^{-\alpha}}{\log u}\,du\right)
 \ll_{\alpha,c}\frac c h\delta
                   \frac{x^{1-\alpha}}{\log x}
 \]
for \(x\) above a threshold depending only on \(c,\alpha\).
This follows from the ratio \((1+\delta)^\alpha=1+O_\alpha(\delta)\);
the last incomplete interval obeys the same bounds.

Write \(A=\sqrt x\,T(x)\) and
\(B=(c/h)x^{1-\alpha}/\log x\).
Taking \(\delta=\min\{1/2,\sqrt{A/B}\}\) gives
\[
 A/\delta+B\delta\ll A+\sqrt{AB}.
 \]
Indeed when \(\sqrt{A/B}\le1/2\) both terms equal \(\sqrt{AB}\);
otherwise \(B<4A\). Since
\[
 \sqrt{AB}
  =\sqrt{c/h}\,x^{3/4-\alpha/2}
       \sqrt{\frac{\log(2x)}{\log x}{\cal L}(x)},
 \]
this proves (A10.1). Its main term is asymptotic to
\(c\,x^{1-\alpha}/[h(1-\alpha)\log x]\).
The ratios of the two error terms to this main term are,
up to fixed constants,
\[
 h{\cal L}(x)(\log x)^2x^{\alpha-1/2},
 \qquad
 \left[h{\cal L}(x)(\log x)^2x^{\alpha-1/2}\right]^{1/2}.
 \]
They tend to zero under (A10.2).
Using the degree-one version of (A09.3) in the identical
partition proves the last assertion.
\end{proof}

\paragraph{NT-A11: a least prime in a prescribed class and arc.}
Let \(0<\mu\le1\) be the length of a fixed arc. Put
\[
 d=\frac h\mu,\qquad
 V=\log(eQ)+n\log(e\,n\,d).
 \]
Under (A09.2), there is an absolute constant \(C>0\) such
that every ray class and every arc of length \(\mu\) contains
a degree-one prime \(\mathfrak p\nmid\mathfrak f\) with
\[
                       N\mathfrak p\le C d^2 V^6 .      \tag{A11.1}
\]
No numerical value of \(C\) is claimed.

\begin{proof}
Use (A09.3) with its degree-one version. Let \(A_0\) be an
absolute error constant there, and put \(x=C d^2 V^6\).
We have \(d\ge1\), \(V\ge n\), \(n\log d\le V\) and
\(n\log V\le2V\). The last inequality follows by writing
\(V=nv\), since \(v\ge\log(e\,n)\) and
\(\log V=\log n+\log v\le2v\).
Consequently, for \(C\ge4\),
\[
 \log x\le(\log C+14)V/n,\qquad
 \log(2x)\le(\log C+15)V/n,\qquad
 {\cal L}(x)\le(32+2\log C)V .
 \]
Also \(\operatorname{Li}_2(x)\ge x/(2\log x)\) for \(x\ge4\).
Thus a sufficient condition for the main term
\(\operatorname{Li}_2(x)/d\) to exceed the error is
\[
 \sqrt C>
 2A_0(\log C+14)(\log C+15)(32+2\log C).
 \]
An absolute \(C\) satisfying this exists, because \(\sqrt C\)
eventually dominates that cubic polynomial in \(\log C\).
The resulting count is positive, proving (A11.1).
\end{proof}

\paragraph{Geometric and run placement.}
The finite condition is a Frobenius fiber in the ray-class
cover of \(\operatorname{Spec}\mathcal O_K[1/\mathfrak f]\);
the angular condition lies in a compact rank-one Hecke torus.
Their joint character algebra is
\(\widehat G\times\mathbf Z\), and (A09.2) is exactly its
Hecke \(L\)-function family. For normalized algebraic Hecke
characters this is the tensor-character description of a CM
torus with a finite Artin factor. For a general unitary
infinity type no algebraic-motive realization is asserted.
The formulas keep the finite mass \(1/h\), the finite
conductor, the infinity parameters and the field degree
separate. They strengthen the separate class and angle
formulations NT-C06/A04--A05. Run:
\texttt{GRH\_NIGHT/number\_theory}; full new proofs are
NT-A08--A11 above, with classical input S-NT-IK and S-NT-ET.
